\documentclass[courseproj]{omgtudoc-asoiu}

\usepackage{amsmath}
\usepackage{booktabs}

\studentname{Морозов}{Алексей}{Иванович}{Морозова}{Алексея}{Ивановича}
\studentgroup{ИВТ-401}
\studentid{260184}
\directionIVT
\courseprojsubject{Методы обработки данных}
\courseprojtheme{Методы проверки качества данных}
\professorvalue{доцент}
\professorinitials{П. П. Петров}
\docyear{2026}

\begin{document}

\maketitle

\begin{abstract}
В этой учебной статье приведён небольшой воспроизводимый пример анализа
качества данных. Структура документа подходит для курсовой работы или
технического отчёта: постановка задачи, методика, результаты и выводы.
\end{abstract}

\section{Постановка задачи}
Перед обработкой набора данных необходимо определить измеряемые показатели,
критерии качества и способ фиксации результатов. В качестве примера будем
рассматривать долю корректно заполненных записей.

\section{Методика}
Пусть $n$ --- число проверенных записей, а $k$ --- число записей без ошибок.
Тогда оценка качества вычисляется по формуле
\begin{equation}
  q = \frac{k}{n}\cdot 100\,\% .
\end{equation}
Порог $q_0$ следует задавать до начала эксперимента, чтобы не подгонять его
под полученный результат.

\begin{table}[ht]
  \centering
  \caption{Пример сводки результатов проверки.}
  \label{tab:quality}
  \begin{tabular}{lrrr}
    \toprule
    Набор & Проверено & Корректно & $q$, \% \\
    \midrule
    A & 120 & 113 & 94.2 \\
    B & 95 & 91 & 95.8 \\
    C & 80 & 72 & 90.0 \\
    \bottomrule
  \end{tabular}
\end{table}

\section{Результаты}
Из таблицы~\ref{tab:quality} видно, что наборы A и B превышают условный
порог в 92\,\%. Для набора C требуется дополнительная проверка причин
пропусков и повторная загрузка исходных данных.

\section{Выводы}
Проверка становится полезной частью рабочего процесса, если её критерии,
исходные данные и вычисления сохраняются вместе с итоговым отчётом. Такой
подход облегчает повторение эксперимента и поиск ошибок.

\end{document}
